\section{Les batteries au lithium}

\setcounter{equation}{0}
\setcounter{Q}{0}

%Composition Chimie A, XEULC, Concours d'admission 2016
Une cellule de batterie lithium-polymère (Li-Po), comme celle qui est représentée sur la Figure 1, peut être constituée d’une électrode de lithium métal, d’un électrolyte contenant des ions \ce{Li+} ($[\ce{Li+}] = 1 ~ \si{\mole\per\liter}$ ) au sein d’un gel polymère organique et d’une électrode composite à base d’oxyde de manganèse de type $\lambda$-\ce{MnO2} (variété cristalline de MnO 2 ). Le composé actif $\lambda$-\ce{MnO2} de structure cubique est synthétisé en plusieurs étapes. La première consiste en la réduction, en solution aqueuse acide, des ions permanganate par l’acide fumarique [acide (E)-but-2-én-1,4-dioïque] pour obtenir l’oxyde de manganèse de type $\beta$-\ce{MnO2} qui a une structure hexagonale.

\Q \'{E}crire la réaction entre les ions \ce{MnO4-} et l’acide fumarique (\ce{C4H4O4}), sachant qu’il n’est pas observé de dégagement de dioxyde de carbone dans ces conditions. Dans une seconde étape, le composé $\beta$-\ce{MnO2} ainsi obtenu est chauffé à \SI{600}{\celsius} sous air pendant plusieurs heures afin d’obtenir la variété $\lambda$-\ce{MnO2} . Cette transformation de structure « hexagonale $\to$ cubique » étant irréversible, commenter la notion de stabilité pour la phase $\beta$-\ce{MnO2}.

Lors de la décharge de la cellule de la batterie (système en fonctionnement), le cristal de $\lambda$-\ce{MnO2} de l’une des électrodes est le siège d’une réaction d’insertion des ions \ce{Li^+}. Cet oxyde de manganèse possède une structure cubique qui peut être décrite comme un réseau cubique à faces centrées (cfc) d’ions \ce{O^{2-}} avec occupation de la moitié des sites octaédriques par des ions \ce{Mn^{4+}} .

\Q Etudier soigneusement la compatibilité de cette description avec la stœchiométrie de l’oxyde.

\Q Ecrire les réactions aux électrodes lors de la décharge de la cellule. Préciser le sens du courant et la polarité des électrodes. Ecrire l’équation-bilan globale lorsque la batterie fonctionne, puis déterminer la tension à vide de la cellule Li-\emph{Po}.

\Q En fait, un seul ion \ce{Mn^{4+}} par maille est susceptible de se réduire en \ce{Mn^{3+}} lors de la décharge de la cellule. Déduire le nombre d’ions \ce{Li+} insérés en site tétraédrique dans cette structure et la formule chimique du composé ainsi obtenu.

\Q La structure initiale $\lambda$-\ce{MnO2} possède un paramètre de maille $a_i = 800 ~ \si{\pico\meter}$ et la structure finale (en fin de décharge de la batterie) un paramètre de maille $a_f = 820 ~ \si{\pico\meter}$. Estimer la variation relative du volume de maille lors de l’insertion de lithium dans la structure de cet oxyde.

\Q Déterminer la capacité spécifique (quantité de charge stockable) exprimée en \si{\ampere\hour\per\kilo\gram} du matériau $\lambda$-\ce{MnO2}. Comparer la valeur estimée avec celle fournie dans le Tableau 1. En déduire la valeur
de l’énergie massique (quantité d’énergie stockable) exprimée en \si{\watt\hour\per\kilo\gram} de ce matériau.

\begin{center}
\begin{tabular}{cccccccc}
 & Chargé & & & & Déchargé & & \\
 & \ce{Li} & \ce{C} & \ce{Si} & \ce{LiMnO2} & \ce{LiC6} & \ce{Li15Si4} & \ce{MnO2} \\
Capacité spécifique [\si{\ampere\hour\per\kilo\gram}] & 3860 & 373 & 3600 & 150 & -- & -- & -- \\ 
Densité & 0.53 & 2.26 & 2.33 & 4.28 & 2.20 & 1.18 & 5.03
\end{tabular}
\captionof{table}{Propriétés de quelques matériaux d'électrode}
\end{center}

Les batteries à base de lithium permettent, dans les conditions d’utilisation recommandées, d’effectuer un nombre élevé de recharges sans perte notable de leur capacité initiale.

\Q \'{E}crire l’équation-bilan de la réaction de recharge de la cellule d’une batterie Li-\emph{Po}. Le redépôt de
lithium métal se fait ordinairement selon un processus de croissance dendritique avec formation de fins filaments qui grandissent perpendiculairement à la surface de l’électrode. Expliquer pour quelle raison cela peut constituer un grave problème de sécurité. Préciser ici l’intérêt d’utiliser un électrolyte gélifié.

Une autre technologie de batterie au lithium permet de remplacer l’électrode métallique de la cellule de l’étude précédente par une électrode en graphite déposée sur un support. Cette technologie est décrite dans le Document 4 qui expose le fonctionnement de la cellule de batterie Li-ion représentée sur le Schéma 3.

\Q Écrire les réactions aux électrodes qui ont lieu au cours du chargement de la batterie Li-ion, puis l’équation-bilan de la réaction. Attribuer l’anode et la cathode. Déterminer la tension à vide d’une cellule de la batterie Li-ion.

Des recherches en cours visent à remplacer l’électrode en graphite de la batterie lithium-ion par une électrode constituée d’un autre élément qui puisse s’allier avec le lithium. Le silicium apparaît de ce point de vue particulièrement prometteur.

\Q \`{A} l’aide du diagramme binaire isobare d’équilibre \ce{Li}-\ce{Si} à $P_{atm}$ de la Figure 2, déterminer la fraction molaire maximale de lithium alliable avec le silicium, à température ambiante.

\Q Le processus de lithiation électrochimique du silicium à \SI{415}{\celsius} est représenté sur la Figure 3. Décrire les étapes successives observées (à y croissant), et interpréter la correspondance entre la courbe de la Figure 3 et le diagramme de la Figure 2.

\Q En réalité, lors de l’insertion électrochimique du lithium dans le silicium à température ambiante, le composé ultime formé est \ce{Li15Si4}. À partir de la position du point représentatif de ce composé dans le diagramme binaire de la Figure 2, préciser la nature des phases solides en présence correspondant à la formule brute \ce{Li15Si4} .

\Q Le silicium cristallise dans un réseau cfc avec occupation d’un site tétraédrique sur deux par des atomes de \ce{Si} (structure diamant). Étudier sans développement excessif si le composé précédent \ce{Li15Si4} est un alliage d’insertion et/ou de substitution.

\Q Connaissant les paramètres de maille du composé initial (\ce{Si}, $a(\ce{Si}) = 540 ~ \si{\pico\meter}$) et final (\ce{Li15Si4}, $a(\ce{Li-Si}) = 1070 ~ \si{\pico\meter}$), estimer la variation relative du volume de maille subi par le silicium. Comparer le comportement des deux électrodes de cette cellule lors de l’insertion du lithium, en s’appuyant sur les résultats de la question 25.

Un constructeur de batteries souhaite comparer les performances et gains obtenus en passant d’un système classique \ce{C} // \ce{Li-MnO2} à un système \ce{Si} // \ce{LiMnO2}.

\Q D’après les données relatives aux matériaux d’électrodes utilisés figurant dans le Tableau 1 et les données commerciales fournies dans le Tableau 2, déterminer la capacité exprimée en \si{\ampere\hour} de chacun des deux systèmes concurrents. Etudier la cohérence des résultats pour chaque système. En déduire l’énergie emmagasinée en Wh par ces deux batteries, et commenter les résultats.

\begin{center}
\begin{tabular}{ccccc}
& Batterie de format 18650 & & Batterie de format 18650 & \\
& \ce{C} // \ce{Li-MnO2} & & \ce{Si} // \ce{Li-MnO2} & \\
& \ce{C} & \ce{LiMnO2} & \ce{Si} & \ce{Li-MnO2} \\
Masse matériau actif [\si{\gram}] & 2.96 & 6.47 & 0.39 & 9.50 \\
Volume batterie [\si{\centi\meter\cubed}] & 16.5 & & 16.5 & \\
\end{tabular}
\captionof{table}{Données commerciales de deux types de batteries Li-ion}
\end{center}

\Q Les batteries contenant une électrode en lithium sont progressivement remplacées par les batteries utilisant la technologie d’insertion des ions lithium. Expliquer si l’étude qui vient d’être décrite conforte ce choix. Etayer la réponse en quelques lignes.

%documents à récuperer dans le sujet + données
